Automation and augmentation do not produce interchangeable model rankings.
Direct-solver performance is relatively concentrated, whereas assistant
performance is more dispersed and task-dependent. The unaided worker winning
three tasks is especially important: guidance can distract the worker,
introduce unnecessary constraints, or redirect it away from the requested
deliverable. Assistant value is therefore conditional, not an automatic
consequence of greater standalone capability.

For model evaluation, the implication is that an automation-only leaderboard
omits a capability relevant to human-in-the-loop and multi-agent deployment.
Benchmarks should distinguish producing a deliverable from improving another
agent's deliverable, because progress on one axis need not transfer to the
other. The fixed-worker design provides a controlled way to add this axis:
holding execution constant makes cross-model differences in guidance
observable without requiring a separate field experiment for every
model--task pair.

The same distinction matters for deployment. A model selected to execute a
task end-to-end may differ from one selected to guide a human, junior model,
or internal workflow. Our experiment uses GPT-3.5-Turbo as a common worker to
identify assistant effects, but the protocol is modular: a deployment-specific
evaluation can substitute the relevant worker while preserving the
cross-assistant comparison. \textsc{CentaurBench} thus complements autonomous
benchmarks and workplace experiments by measuring role-specific performance
under a shared, reproducible protocol.
